\chapter{Đứng Dậy Sau Vấp Ngã}
\markboth{Đứng Dậy Sau Vấp Ngã}{Rising After a Fall}

\section{Ngã Không Đáng Xấu Hổ, Nằm Luôn Mới Đáng Lo}
\begin{truyen}{Ngã Không Đáng Xấu Hổ, Nằm Luôn Mới Đáng Lo}{Falling Is Not Shameful, Staying Down Is the Real Danger}
\chuhoa{M}{ột học sinh thi trượt, bị chê cười, rồi tự kết luận đời mình kém cỏi.}
\textit{(A student fails an exam, gets mocked, and concludes that life itself is inferior.)}

Nhưng cuộc đời không đo con người bằng số lần ngã.
\textit{(But life does not measure a person by the number of falls.)}

Nó đo bằng chuyện sau đó em làm gì.
\textit{(It measures what happens next.)}

Có người vấp rồi học được cách đi chắc hơn.
\textit{(Some people stumble and learn to walk more firmly.)}

Có người chỉ mải xấu hổ nên không đứng dậy nữa.
\textit{(Others stay busy feeling ashamed and never rise again.)}

Vấp ngã là chuyện thường.
\textit{(Falling is ordinary.)}

Bỏ mặc bản thân sau cú ngã mới là điều đáng sợ.
\textit{(Abandoning yourself after the fall is what is frightening.)}

\end{truyen}

\begin{baihoc}
Điều cần dạy học sinh không phải là làm sao không bao giờ ngã. Điều cần dạy là cách tự kéo mình dậy sau khi ngã.
\textit{(What students need is not a way never to fall. They need a way to lift themselves after falling.)}
\end{baihoc}

\begin{ghinhoanh}
\textbf{Short English to remember:}

What students need is not a way never to fall.

They need a way to lift themselves after falling.

\end{ghinhoanh}

\begin{tuhocnhanh}
\textbf{Cau hoi tu hoc / Self-study questions}

1. Van de chinh la gi? \textit{What is the main problem?}

2. Cach xu ly tot hon la gi? \textit{What is a better action?}

3. Mot cau tieng Anh em muon nho la cau nao? \textit{Which English sentence do you want to remember?}
\end{tuhocnhanh}
\ngancach

\section{Bắt Đầu Lại Luôn Nhỏ Và Xấu}
\begin{truyen}{Bắt Đầu Lại Luôn Nhỏ Và Xấu}{Starting Again Is Always Small and Ugly}
\chuhoa{N}{hiều em tưởng đứng dậy phải thật hoành tráng: lập kế hoạch lớn, thay đổi triệt để, hứa với bản thân từ nay khác hẳn.}
\textit{(Many students imagine recovery must be dramatic: a huge plan, total transformation, promises of becoming completely different from now on.)}

Sự thật là bắt đầu lại thường rất nhỏ và khá xấu.
\textit{(The truth is that starting again is usually very small and rather ugly.)}

Một buổi ngồi vào bàn học lại.
\textit{(One evening sitting at the desk again.)}

Một bài tập làm cho xong.
\textit{(One assignment completed.)}

Một cuộc nói chuyện trung thực.
\textit{(One honest conversation.)}

Một giấc ngủ đúng giờ.
\textit{(One proper night of sleep.)}

Những thứ nhỏ ấy không gây xúc động mạnh, nhưng chúng mới thật sự dựng lại cuộc đời.
\textit{(Those small things do not create dramatic emotion, but they are what actually rebuild a life.)}

Người thích chờ động lực lớn thường chờ rất lâu.
\textit{(People who wait for grand motivation often wait a very long time.)}

Người chấp nhận bước nhỏ thì đi được.
\textit{(Those who accept small steps actually move.)}

\end{truyen}

\begin{baihoc}
Dạy học sinh yêu bước nhỏ. Hồi phục bền gần như luôn bắt đầu bằng thứ rất bình thường.
\textit{(Teach students to respect small steps. Durable recovery almost always begins with something ordinary.)}
\end{baihoc}

\begin{ghinhoanh}
\textbf{Short English to remember:}

Teach students to respect small steps.

Durable recovery almost always begins with something ordinary.

\end{ghinhoanh}

\begin{tuhocnhanh}
\textbf{Cau hoi tu hoc / Self-study questions}

1. Van de chinh la gi? \textit{What is the main problem?}

2. Cach xu ly tot hon la gi? \textit{What is a better action?}

3. Mot cau tieng Anh em muon nho la cau nao? \textit{Which English sentence do you want to remember?}
\end{tuhocnhanh}
\ngancach

\section{Đời Không Công Bằng, Nhưng Em Vẫn Có Thể Chơi Đẹp}
\begin{truyen}{Đời Không Công Bằng, Nhưng Em Vẫn Có Thể Chơi Đẹp}{Life Is Not Fair, But You Can Still Play Clean}
\chuhoa{C}{ó người sinh ra đã có điều kiện hơn, nền tốt hơn, đỡ áp lực hơn.}
\textit{(Some people are born with better conditions, stronger support, and less pressure.)}

Điều đó là thật.
\textit{(That is real.)}

Học sinh cần được nói sự thật này, thay vì nghe những lời màu hồng rằng ai cố cũng như ai.
\textit{(Students need to hear this truth rather than soft slogans claiming that everyone’s effort starts equally.)}

Nhưng từ bất công đó không suy ra rằng em được quyền buông xuôi, làm bậy, hay dối trá.
\textit{(But that unfairness does not imply that you have the right to give up, act wrongly, or become dishonest.)}

Đời không công bằng là lý do để em tỉnh táo hơn, không phải bẩn hơn.
\textit{(Life being unfair is a reason to become clearer, not dirtier.)}

\end{truyen}

\begin{baihoc}
Giáo dục trung thực phải nhìn nhận chênh lệch thật của đời sống, nhưng vẫn giữ cho học sinh một nguyên tắc cứng: khó không cho phép mình trở nên xấu.
\textit{(Honest education must acknowledge real inequality in life, while preserving one hard principle: difficulty does not permit becoming corrupt.)}
\end{baihoc}

\begin{ghinhoanh}
\textbf{Short English to remember:}

Honest education must acknowledge real inequality in life, while preserving one hard principle: difficulty does not permit becoming corrupt.

\end{ghinhoanh}

\begin{tuhocnhanh}
\textbf{Cau hoi tu hoc / Self-study questions}

1. Van de chinh la gi? \textit{What is the main problem?}

2. Cach xu ly tot hon la gi? \textit{What is a better action?}

3. Mot cau tieng Anh em muon nho la cau nao? \textit{Which English sentence do you want to remember?}
\end{tuhocnhanh}
\ngancach

\section{Cuối Cùng, Điều Giữ Em Đi Tiếp Không Phải Là Lời Khen}
\begin{truyen}{Cuối Cùng, Điều Giữ Em Đi Tiếp Không Phải Là Lời Khen}{In the End, Praise Is Not What Keeps You Going}
\chuhoa{C}{ó thời gian học sinh sống bằng lời khen. Được công nhận là có năng lượng. Bị chê là rơi xuống ngay.}
\textit{(For a time, students may live on praise. Recognition gives energy; criticism drops them instantly.)}

Nhưng càng lớn, em càng phải đi bằng thứ bền hơn: giá trị bên trong, thói quen lao động, và khả năng tự nhắc mình lý do phải tiếp tục.
\textit{(But as you grow, you must move by something more durable: inner values, habits of work, and the ability to remind yourself why to continue.)}

Không ai được vỗ tay cho em mọi ngày.
\textit{(No one will applaud you every day.)}

Có những giai đoạn rất dài em phải bước trong im lặng.
\textit{(There are long phases when you must walk in silence.)}

Người đi đường dài là người không chết đói vì thiếu lời khen.
\textit{(The person who lasts is the one who does not starve in the absence of praise.)}

\end{truyen}

\begin{baihoc}
Hãy giúp học sinh xây nội lực thay vì chỉ bơm cảm xúc. Người trưởng thành phải đi được cả khi không ai cổ vũ.
\textit{(Help students build inner strength rather than only emotional boosts. A mature person must be able to continue even when no one cheers.)}
\end{baihoc}

\begin{ghinhoanh}
\textbf{Short English to remember:}

Help students build inner strength rather than only emotional boosts.

A mature person must be able to continue even when no one cheers.

\end{ghinhoanh}

\begin{tuhocnhanh}
\textbf{Cau hoi tu hoc / Self-study questions}

1. Van de chinh la gi? \textit{What is the main problem?}

2. Cach xu ly tot hon la gi? \textit{What is a better action?}

3. Mot cau tieng Anh em muon nho la cau nao? \textit{Which English sentence do you want to remember?}
\end{tuhocnhanh}
\ngancach

\section{Đứng Dậy Không Có Nhạc Nền}
\begin{truyen}{Đứng Dậy Không Có Nhạc Nền}{Rising Without Background Music}
\chuhoa{N}{hiều câu chuyện truyền cảm hứng kể việc đứng dậy như một khoảnh khắc rất đẹp.}
\textit{(Many inspirational stories portray rising again as a beautiful moment.)}

Đời thật thường không vậy.
\textit{(Real life usually does not.)}

Đứng dậy ngoài đời thường diễn ra trong căn phòng lộn xộn, tâm trạng nặng, và một gương mặt chưa hề thấy khá hơn.
\textit{(Rising in real life often happens in a messy room, with a heavy mood, and a face that does not yet look better.)}

Em vẫn buồn, vẫn ngại, vẫn chưa tin mình ổn lên được, nhưng em làm việc cần làm trước đã.
\textit{(You may still feel sad, embarrassed, and unconvinced that things will improve, but you do the necessary work first.)}

Hồi phục thật thường không kịch tính.
\textit{(Real recovery is often not dramatic.)}

Nó bền vì nó không phụ thuộc cảm hứng.
\textit{(It lasts because it does not depend on inspiration.)}

\end{truyen}

\begin{baihoc}
Đừng chờ cảm giác sẵn sàng mới đứng dậy. Nhiều lần, hành động đi trước rồi tinh thần mới theo sau.
\textit{(Do not wait to feel ready before rising. Many times action comes first, and the spirit follows later.)}
\end{baihoc}

\begin{ghinhoanh}
\textbf{Short English to remember:}

Do not wait to feel ready before rising.

Many times action comes first, and the spirit follows later.

\end{ghinhoanh}

\begin{tuhocnhanh}
\textbf{Cau hoi tu hoc / Self-study questions}

1. Van de chinh la gi? \textit{What is the main problem?}

2. Cach xu ly tot hon la gi? \textit{What is a better action?}

3. Mot cau tieng Anh em muon nho la cau nao? \textit{Which English sentence do you want to remember?}
\end{tuhocnhanh}
\ngancach

\section{Người Thợ Học Lại Từ Mũi Đục Đầu Tiên}
\begin{truyen}{Người Thợ Học Lại Từ Mũi Đục Đầu Tiên}{The Craftsman Who Started Again With the First Chisel Stroke}
\chuhoa{C}{ó người thợ làm hỏng một món lớn, mất uy tín, phải bắt đầu lại từ những việc nhỏ nhất.}
\textit{(A craftsman ruined an important commission, lost credibility, and had to start again from the smallest jobs.)}

Ông không lấy lại tên tuổi bằng một cú bứt phá.
\textit{(He did not regain his name through one dramatic comeback.)}

Ông lấy lại nó bằng hàng trăm lần làm đúng những việc bé mà không ai vỗ tay.
\textit{(He regained it through hundreds of correct small acts that no one applauded.)}

Học sinh sau cú ngã cũng vậy.
\textit{(Students after a fall are the same.)}

Danh dự và niềm tin thường được xây lại bằng việc nhỏ lặp đi lặp lại.
\textit{(Dignity and trust are often rebuilt through repeated small actions.)}

Không có đường tắt tử tế cho việc dựng lại một con người.
\textit{(There is no decent shortcut for rebuilding a person.)}

\end{truyen}

\begin{baihoc}
Muốn đứng dậy bền, hãy chấp nhận làm lại từ những viên gạch nhỏ. Sự nóng vội thường chỉ làm căn nền mới nứt tiếp.
\textit{(To rise in a lasting way, accept rebuilding with small bricks. Impatience often only cracks the new foundation again.)}
\end{baihoc}

\begin{ghinhoanh}
\textbf{Short English to remember:}

To rise in a lasting way, accept rebuilding with small bricks.

Impatience often only cracks the new foundation again.

\end{ghinhoanh}

\begin{tuhocnhanh}
\textbf{Cau hoi tu hoc / Self-study questions}

1. Van de chinh la gi? \textit{What is the main problem?}

2. Cach xu ly tot hon la gi? \textit{What is a better action?}

3. Mot cau tieng Anh em muon nho la cau nao? \textit{Which English sentence do you want to remember?}
\end{tuhocnhanh}
\ngancach

\section{Bài Học Không Có Trong Bằng Khen}
\begin{truyen}{Bài Học Không Có Trong Bằng Khen}{The Lesson Not Found in Certificates}
\chuhoa{B}{ằng khen ghi lại lúc em đang sáng.}
\textit{(Certificates record moments when you shine.)}

Nhưng phần trưởng thành lớn nhất nhiều khi được viết ở những đoạn tối hơn.
\textit{(But the greatest part of maturity is often written in darker passages.)}

Một người học cách ngồi dậy sau trượt ngã sẽ mang kỹ năng đó suốt đời, trong công việc, hôn nhân, bệnh tật, và mất mát.
\textit{(A person who learns to rise after falling carries that skill for life, into work, marriage, illness, and loss.)}

Không có tấm giấy đẹp nào ghi: em đã không bỏ mình lại sau khi thất bại.
\textit{(There is no elegant certificate that says: you did not abandon yourself after failure.)}

Nhưng đó là một thành tựu rất lớn.
\textit{(Yet that is a great achievement.)}

Trường học nếu chỉ khen thành tích mà không dạy hồi phục thì đang bỏ quên một nửa đời người.
\textit{(A school that praises achievement but does not teach recovery is neglecting half of life.)}

\end{truyen}

\begin{baihoc}
Giáo dục cần tôn trọng những kỹ năng không treo trên tường: hồi phục, bền bỉ, và tự kéo mình qua giai đoạn tệ.
\textit{(Education needs to honor skills that are never hung on walls: recovery, endurance, and self-restoration through hard periods.)}
\end{baihoc}

\begin{ghinhoanh}
\textbf{Short English to remember:}

Education needs to honor skills that are never hung on walls: recovery, endurance, and self-restoration through hard periods.

\end{ghinhoanh}

\begin{tuhocnhanh}
\textbf{Cau hoi tu hoc / Self-study questions}

1. Van de chinh la gi? \textit{What is the main problem?}

2. Cach xu ly tot hon la gi? \textit{What is a better action?}

3. Mot cau tieng Anh em muon nho la cau nao? \textit{Which English sentence do you want to remember?}
\end{tuhocnhanh}
\ngancach

\section{Không Cần Trả Thù Thất Bại}
\begin{truyen}{Không Cần Trả Thù Thất Bại}{You Do Not Need to Take Revenge on Failure}
\chuhoa{S}{au khi thua, nhiều học sinh lao vào học như để trả thù cú ngã vừa rồi.}
\textit{(After losing, many students throw themselves into study as if taking revenge on the fall.)}

Năng lượng đó có thể mạnh trong ngắn hạn, nhưng thường mang theo rất nhiều cay cú và tự ép.
\textit{(That energy may be powerful in the short term, but it often carries bitterness and self-violence.)}

Đứng dậy không nhất thiết phải dựa vào thù hằn.
\textit{(Rising again does not have to rely on resentment.)}

Em có thể đi tiếp bằng sự điềm tĩnh hơn: mình sai đâu sửa đó, chưa đủ đâu bồi đó.
\textit{(You can continue more calmly: fix where you were wrong, strengthen where you were not enough.)}

Người phục hồi lành mạnh thường đi xa hơn người chỉ cháy bùng lên vì tức.
\textit{(Those who recover healthily often go farther than those who burn briefly from anger.)}

\end{truyen}

\begin{baihoc}
Đừng biến thất bại thành cuộc chiến với chính mình. Sửa lại bản thân hiệu quả hơn trừng phạt bản thân.
\textit{(Do not turn failure into war against yourself. Correcting yourself works better than punishing yourself.)}
\end{baihoc}

\begin{ghinhoanh}
\textbf{Short English to remember:}

Do not turn failure into war against yourself.

Correcting yourself works better than punishing yourself.

\end{ghinhoanh}

\begin{tuhocnhanh}
\textbf{Cau hoi tu hoc / Self-study questions}

1. Van de chinh la gi? \textit{What is the main problem?}

2. Cach xu ly tot hon la gi? \textit{What is a better action?}

3. Mot cau tieng Anh em muon nho la cau nao? \textit{Which English sentence do you want to remember?}
\end{tuhocnhanh}
\ngancach

\section{Đường Dài Ưa Người Bền Hơn Người Bốc}
\begin{truyen}{Đường Dài Ưa Người Bền Hơn Người Bốc}{The Long Road Favors the Steady Over the Explosive}
\chuhoa{C}{ó học sinh sau một cú vấp bùng lên rất mạnh vài tuần rồi tắt.}
\textit{(Some students erupt with effort for a few weeks after a fall and then burn out.)}

Có học sinh khác hồi lại chậm hơn nhưng đều.
\textit{(Others recover more slowly but steadily.)}

Đường dài hiếm khi thưởng cho người bốc nhất.
\textit{(The long road rarely rewards the most explosive person.)}

Nó thường thưởng cho người giữ nhịp được lâu.
\textit{(It usually rewards the one who can keep rhythm.)}

Đời sống học tập cũng như tập luyện cơ thể: quá đà rồi đứt quãng thường kém hơn vừa sức mà bền.
\textit{(Study is like training the body: going too hard and breaking off is often worse than a moderate pace sustained well.)}

Sau vấp ngã, việc cần nhất không phải là hứa thật lớn.
\textit{(After a stumble, the most necessary thing is not a grand promise.)}

Nó là tìm lại nhịp có thể giữ được.
\textit{(It is recovering a rhythm you can sustain.)}

\end{truyen}

\begin{baihoc}
Muốn đứng dậy và đi xa, học sinh phải học yêu sự bền hơn sự bốc đồng. Nhịp đúng quan trọng hơn hào quang ngắn.
\textit{(To rise and travel far, students must learn to value steadiness over sudden intensity. The right rhythm matters more than brief glory.)}
\end{baihoc}

\begin{ghinhoanh}
\textbf{Short English to remember:}

To rise and travel far, students must learn to value steadiness over sudden intensity.

The right rhythm matters more than brief glory.

\end{ghinhoanh}

\begin{tuhocnhanh}
\textbf{Cau hoi tu hoc / Self-study questions}

1. Van de chinh la gi? \textit{What is the main problem?}

2. Cach xu ly tot hon la gi? \textit{What is a better action?}

3. Mot cau tieng Anh em muon nho la cau nao? \textit{Which English sentence do you want to remember?}
\end{tuhocnhanh}
\ngancach

\section{Một Ngày Nhìn Lại, Em Sẽ Biết Mình Đã Qua}
\begin{truyen}{Một Ngày Nhìn Lại, Em Sẽ Biết Mình Đã Qua}{One Day You Will Look Back and Know You Passed Through}
\chuhoa{K}{hi đang ở giữa giai đoạn tệ, học sinh rất khó tin rằng rồi nó sẽ lùi lại phía sau.}
\textit{(When in the middle of a bad period, students find it very hard to believe it will eventually move behind them.)}

Nhưng nhiều người trưởng thành nhìn lại đều có một đoạn từng tưởng không qua nổi, cuối cùng vẫn qua.
\textit{(Yet many adults look back and find a period they once thought unsurvivable, and still they passed through it.)}

Không phải vì lúc đó họ mạnh sẵn.
\textit{(Not because they were already strong then.)}

Nhiều khi chỉ vì họ chịu sống thêm từng ngày một, rồi từng bước một.
\textit{(Often simply because they kept living one day at a time, and then one step at a time.)}

Hồi phục lớn đôi khi chỉ là tổng cộng của rất nhiều ngày không bỏ cuộc.
\textit{(Great recovery is sometimes nothing more than the sum of many days of not quitting.)}

\end{truyen}

\begin{baihoc}
Khi học sinh đang ở đáy, đừng bắt các em nhìn quá xa. Nhiều lúc mục tiêu đủ tốt chỉ là qua được hôm nay cho tử tế.
\textit{(When students are at the bottom, do not force them to look too far ahead. Sometimes the goal that is good enough is simply to get through today with dignity.)}
\end{baihoc}

\begin{ghinhoanh}
\textbf{Short English to remember:}

When students are at the bottom, do not force them to look too far ahead.

Sometimes the goal that is good enough is simply to get through today with dignity.

\end{ghinhoanh}

\begin{tuhocnhanh}
\textbf{Cau hoi tu hoc / Self-study questions}

1. Van de chinh la gi? \textit{What is the main problem?}

2. Cach xu ly tot hon la gi? \textit{What is a better action?}

3. Mot cau tieng Anh em muon nho la cau nao? \textit{Which English sentence do you want to remember?}
\end{tuhocnhanh}
